\chapter{Reference Guide and Theorem Crosswalk}
\label{app:perverse-references}

The notes are designed to be read linearly; this appendix identifies the
original locations of the deep results and records how their conventions map
to ours.

\section{BBD}

The primary source for the categorical language is Beilinson--Bernstein--
Deligne, \emph{Faisceaux pervers} \cite{BBD}.

\begin{longtable}{
 >{\raggedright\arraybackslash}p{0.24\textwidth}
 >{\raggedright\arraybackslash}p{0.20\textwidth}
 >{\raggedright\arraybackslash}p{0.46\textwidth}}
\toprule
Topic & BBD location & Use here\\
\midrule
\endhead
Middle perversity and gluing &
Sections 1.4, 2.1--2.2 &
\cref{def:middle-perversity,thm:middle-t-structure}.\\

Affine exactness &
Theorem 4.1.1 and Corollaries 4.1.2--4.1.3 &
\cref{prop:basic-perverse-exactness,cor:perverse-artin-vanishing}.\\

Finite length and simple objects &
Theorem 4.3.1 &
\cref{thm:perverse-finite-length,thm:simple-perverse-classification}.\\

Intermediate extension and weights &
Corollary 5.3.2 &
\cref{thm:jmiddle-purity}; used only after the direct-image theorem.\\

Pure perverse semisimplicity &
Theorem 5.3.8 &
\cref{thm:pure-perverse-geometric-semisimple}.\\

Ordinary/perverse weight criterion &
Theorem 5.4.1 and Corollary 5.4.4 &
\cref{thm:perverse-weight-criterion}.\\

Splitting of a pure complex &
Theorem 5.4.5 &
\cref{thm:pure-complex-splitting}.\\

Decomposition theorem &
Theorem 6.2.5 &
\cref{thm:decomposition-finite-field} and the characteristic-zero
spreading-out discussion.\\
\bottomrule
\end{longtable}

BBD's weight chapters use Deligne's \emph{Weil II}.  Therefore the entries in
Sections 5--6 of BBD are consequences, not independent inputs, in this book.

\section{Deligne's Weil II}

Deligne's \emph{La conjecture de Weil II} \cite{DeligneWeilII} supplies the
arithmetic weight theory.

\begin{longtable}{
 >{\raggedright\arraybackslash}p{0.27\textwidth}
 >{\raggedright\arraybackslash}p{0.18\textwidth}
 >{\raggedright\arraybackslash}p{0.45\textwidth}}
\toprule
Result & Location & Role\\
\midrule
\endhead
Weights and determinant weights &
Sections 1.2--1.3 &
Conventions and monodromy bookkeeping in \cref{chap:weights}.\\

Real-sheaf criterion &
Theorem 1.5.1 &
Mixedness input developed in the real-sheaf chapter.\\

Local monodromy weights &
Theorem 1.8.4 &
Boundary control in the curve estimate.\\

Curve purity &
Theorem 3.2.3 &
The terminal result reconstructed by the Fourier chapters.\\

General direct-image upper bound &
Theorem 3.3.1 &
\cref{thm:deligne-laumon-direct-image}.\\

Compact-support and dual consequences &
Corollaries 3.3.4--3.3.6 &
\cref{cor:perverse-upper-direct-image,cor:perverse-lower-direct-image}.\\

Smooth-proper consequences &
Corollary 3.3.9 &
Recovered in
\cref{cor:perverse-smooth-proper-purity,chap:weil-assembly,chap:integrality}.\\
\bottomrule
\end{longtable}

\section{Laumon's Fourier transform}

Laumon's \emph{Transformation de Fourier, constantes d'equations
fonctionnelles et conjecture de Weil} \cite{LaumonFourier} gives the
perverse-Fourier route used here.

\begin{longtable}{
 >{\raggedright\arraybackslash}p{0.26\textwidth}
 >{\raggedright\arraybackslash}p{0.20\textwidth}
 >{\raggedright\arraybackslash}p{0.44\textwidth}}
\toprule
Result & Location & Convention/use\\
\midrule
\endhead
Definition &
1.2.1.1 &
\(\FT_\psi(K)=Rp_{2!}(p_1^*K\otimes\cL_\psi(\langle,\rangle))[r]\).\\

Inversion &
Theorem 1.2.2.1 &
With inverse character,
\(\FT_{\psi^{-1}}\FT_\psi(K)\simeq K(-r)\); with the same character,
insert pullback by \([-1]\).\\

Forget-support isomorphism &
Theorem 1.3.1.1 &
For the Fourier kernel, the relevant \(Rp_!\to Rp_*\) map is an isomorphism.\\

Duality &
Corollary 1.3.2.2 &
The character changes from \(\psi\) to \(\psi^{-1}\), with the stated twist.\\

Perverse \(t\)-exactness &
Theorem 1.3.2.3 and Corollary 1.3.2.4 &
Exact equivalence of perverse hearts and preservation of simple objects.\\

Local Fourier transforms &
Sections 2.3--2.5 &
Local monodromy and boundary exact sequences.\\

Curve weight proof &
Sections 4.1--4.4 &
The curve-estimate chapter.\\
\bottomrule
\end{longtable}

\section{Expository guides}

Kiehl--Weissauer \cite{KiehlWeissauer} gives the closest book-length
presentation of the same proof technology.  The Stanford seminar
\cite{StanfordWeilIISeminar}, especially Lectures 18--24
\cite{StanfordL18,StanfordL19,StanfordL22,StanfordL23,StanfordL24}, is an
excellent route map.  Its lecture notes are informal; where wording or
normalization is delicate, the theorem statements in Deligne, Laumon, and BBD
govern this text.

\section{Convention crosswalk}

\begin{longtable}{
 >{\raggedright\arraybackslash}p{0.30\textwidth}
 >{\raggedright\arraybackslash}p{0.29\textwidth}
 >{\raggedright\arraybackslash}p{0.31\textwidth}}
\toprule
Choice in these notes & Consequence & Conversion warning\\
\midrule
\endhead
Geometric Frobenius &
\(F|\Qbarl(1)=q^{-1}\). &
Arithmetic Frobenius inverts all eigenvalues.\\

Perverse-normalized IC &
\(\IC_X(\cL)=j_{!*}\cL[d]\). &
Sources using unshifted IC differ by \([-d]\).\\

Pure complex convention &
\(\cH^aK\) has weight \(w+a\). &
Thus \(K[n]\) has weight \(w+n\).\\

Fourier shift by rank &
\(\FT_\psi[-]\) includes \([r]\). &
Without it, Fourier is not perverse \(t\)-exact.\\

Frobenius polynomial &
\(P_i(T)=\det(1-TF)\). &
Its polynomial roots are reciprocal to the Frobenius eigenvalues.\\
\bottomrule
\end{longtable}

These five choices determine every sign, shift, twist, and absolute value in
the final theorem.
